\paragraph{Degrees of the CM lines.}

The approximation determines the CM lines. Their contribution to the section height involves both the finite ideal and the norms of its chosen frame at infinity.

For a metrized line \(L\) over a number field let
\(I(L)\) be its fractional ideal and \(q_\sigma(L)\) its
metric on the chosen coordinate at \(\sigma\). We use the
normalized degree
\[
 \operatorname{ndeg}L
   =-\frac{\sum_w\operatorname{ord}_w I(L)\log\mathrm Nw+
                   \sum_\sigma\log q_\sigma(L)}
              {[K':\mathbb Q]} .
\]
The product formula makes this independent of the coordinate.
For the CM object of dimension \(g\), use its Riemann
metric on the Lie lattice \(\mathrm{Lie}\) and the integral
theta line \(\mathcal T\) with the normalized \(L^2\) metric.
\begin{lemma}[The principal cubical theta degree]\label{lem:slice-principal-theta-degree}
For this CM Lie lattice and the normalized principal theta line,
the degree comparison is
\begin{equation}\label{eq:slice-cm-theta-degree}
 \operatorname{ndeg}(\mathcal T^\vee)
      =-\tfrac12\operatorname{ndeg}(\det\mathrm{Lie})
                         +\tfrac{g}{4}\log(2\pi).
\end{equation}
\end{lemma}

\begin{proof}
The theta generator comes from the rigidification of
the principal symmetric cubical line. The analytic length
is the mean of its squared Gaussian norm on a period
parallelepiped, with its square root taken, and the Lie
determinant length is the square root of the Riemann
determinant in the same frame. To see the comparison, write
\(\mathcal D=\det\mathrm{Lie}\) and use the powered cubical invariant
of the rigidified symmetric principal line. For the chosen theta
and Lie coordinates it is a scalar \(k\in (K')^\times\) for an
integer power parameter \(a>0\) depending on the dimension.
Multiplying the theta generator by \(u\) and changing the Lie
frame by \(M\) multiplies \(k\) by \(u^{-8a}\det(M)^{4a}\).
The invariant in integral theta and Lie generators is a unit, by
the integral cubical identity.

Comparing those generators with
the chosen coordinates gives
\[
 (k)=I(\mathcal T)^{8a}I(\mathcal D)^{-4a}.
\]
This is an equality of fractional ideals. One may compare over a
valued extension, since both sides acquire the same factor of the
ramification index. At a complex embedding \(\sigma\), the same
invariant in a theta uniformization satisfies
\[
 |\sigma(k)|\,q_\sigma(\mathcal T)^{8a}
       q_\sigma(\mathcal D)^{-4a}=(2\pi)^{2ga}.
\]
Here the complex cubical formula uses the mean square of the theta
generator, and the period quotient has volume
\(\pi^g |\det W|^2/q_\sigma(\mathcal D)^2\) in the corresponding
complex uniformization frame \(W\). Thus the volume factor uses
the Lie determinant of the same frame. Let \(d=[K':\mathbb Q]\). Summing the complex identities and
using the product formula for \(k\) replaces
\(\sum_\sigma\log|\sigma(k)|\) by the finite log norm of the
displayed ideal. Thus
\[
      -8ad\,\operatorname{ndeg}(\mathcal T)
             +4ad\,\operatorname{ndeg}(\mathcal D)
             =2gad\log(2\pi).
\]
Divide by \(8ad\) and use duality for the degree of a
metrized line. \label{cor:slice-integral-norm-1}
\end{proof}

For an abelian CM differential line \(\ell_k\), put \(\Omega=2\operatorname{ndeg}\ell_k\). The CM differential degree formula gives the same real number for every abelian line in both constructions. The toric differential lines on the arc have degree zero.
\label{cor:slice-integral-norm-2}

\begin{lemma}[Degree of a CM coefficient line]\label{lem:cm-coefficient-line-degree}
With the theta degree identity \eqref{eq:slice-cm-theta-degree}
and the abelian letter metrics above, the value line of an arc
coefficient satisfies \eqref{eq:slice-cm-line-degree}. For the coefficient line \(\mathcal L_\mu\), the bound is
\[
 \operatorname{ndeg}\mathcal L_\mu\le P_\Lambda(\mathcal A+frn),\qquad P_\Lambda=\tfrac12(1+|\Omega|+\log\Lambda).
\]

\end{lemma}

\begin{proof}
\proofstep{Compensating the norms of the CM letters.}
The abelian letters \(u_k\) form a basis over \(K'\). Each
scaled line \(I(\ell_k)u_k\) is inside the completed Lie
lattice. Thus their exterior product is inside its determinant
lattice. Over each complex embedding their lengths are at
most \(\Lambda q_\sigma(\ell_k)\).
To combine the two estimates, identify the exterior line of
the \(u_k\) with \(K'\) via that exterior product. The
determinant ideal contains \(\prod_k I(\ell_k)\), and
Hadamard bounds the determinant metric by
\(\Lambda^g\prod_k q_\sigma(\ell_k)\).

Taking degrees with the indicated sign gives
\[
 \operatorname{ndeg}(\det\mathrm{Lie})
             \ge g\,\Omega/2-g\log\Lambda .
\]
By \eqref{eq:slice-cm-theta-degree},
\(\operatorname{ndeg}(\mathcal T^\vee)\le g P_\Lambda\),
where
\[
             P_\Lambda=\frac{1+|\Omega|}{2}
                                  +\tfrac12\log\Lambda .
\]
Here \(\log(2\pi)<2\). Every dual letter has degree
at most \(P_\Lambda\) since \(\Lambda\ge1\).
The arc has abelian dimension \(f(r-2)\le fr\).
The line \(\mathcal L_\mu\) of the arc coefficient therefore
has
\begin{equation}\label{eq:slice-cm-line-degree}
    \operatorname{ndeg}\mathcal L_\mu
           \le P_\Lambda(\mathcal A+frn).
\end{equation}
This is exactly the degree of the line in the finite and
complex comparisons. All three bounds for the arc
coefficient now have their dependence on \(\Lambda\), \(n'\),
\(\Omega\), and the actual cutoff specified.
\label{cor:slice-integral-norm-3}

\end{proof}

\paragraph{Uniformity of the radius and torsion order.}

The detecting CM point may depend on the form. The approximation radius and the allowed torsion order must nevertheless be bounded before that point is chosen.

Put \(w_0=S_\kappa+r\), \(N_A=d_E=|\operatorname{disc}E|\),
\(N_C=\mathrm N\mathfrak c_K\), and \(M_0=g_0g_1\).
The split approximation compares the two cusp times using
a positive time lower bound \(e_{\rm sp}\), and uses
\(\rho_{\rm sp}=1/(12\pi(c_1+w_0))\) for the approximation.
With
\begin{align*}
 N_0&=M_0d_E^2,\qquad N_2=4f,\\
 H_0&=N_C^2(2+4f^4\operatorname{rd}(E)^2 N_A^2),\\
 K_0'&=1+N_0H_0(2c_1+1),\\
 R_0&=w_0+N_2\log(2K_0' N_2)+N_2+2 ,
\end{align*}
the time bound is
\[
 e_{\rm sp}=\min\{1,f^2/(N_0^2N_C^2c_1R_0)\},
 \qquad e_{\rm sp}^{-1}\le\max\{1,N_0^2N_C^2c_1R_0\}.
\] The coefficients in the bounds
for the split complementary frame are
\begin{align*}
 B_{\rm ind}&=f(2f(r-2))^3 d_E N_A
                     \sqrt{N_C^2\operatorname{rd}(E)e_{\rm sp}^{-1}},\\
 B_{\rm rad}&=2\cdot4\sqrt{r-2}\cdot
         4(f(r-2))^3d_E\sqrt{e_{\rm sp}^{-1}}\cdot40\rho_{\rm sp}^{-1}.
\end{align*}

The choice is \(n'=1+\lceil1000B_{\rm ind}\rho_{\rm sp}^{-1}\rceil\)
and \(\Lambda=\max(1,B_{\rm rad}^{2q_j-1})\).
Fix this order before choosing the detecting frame.
The split approximation theorem then chooses a frame
where multiplication by this order makes the abelian
projection of the cusp lattice integral. The constants
are the same for all choices of arc on a lift.

\begin{lemma}[Uniform order and radius for the detecting arcs]\label{lem:uniform-cm-arc-size}
Use the split approximation with the selected weight and
cutoff \(c_1\), and the fixed parameters above. For \(z=frL\),
the order \(n'\) and radius bound \(\Lambda\) satisfy
\[
 \max(\log n',\log\Lambda)\le200(K_*z)^2,\qquad K_*=20+8|C_A|+|C_K|.
\]
 These choices precede selection
of a detecting arc. The three coefficient allowances in
the proof each have the bound
\((17+2K_*+200K_*^4)z^4S_\kappa\).
\end{lemma}

\begin{proof}
\proofstep{Uniform bounds for the conductor and radius.}
Here is the quantitative consequence used in the division
budget. Use \(z=frL\) as in the run and
\(K_*=20+|C_A|+7|C_A|+|C_K|\).
The size inequalities put \(f,r,d_E,M_0,N_C,c_1,w_0\)
below 
\[
  U=e^{K_*z}.
\]
 For \(c_1\) this uses
\(Tc\le1\) and \(q_j\le r\). For \(w_0\) use
\(S_\kappa\le C_Kkfb\), \(k\le r\), \(\log b=L\).
The size comparison also gives
\(|\Omega|\le |C_A|L\le K_*z\).
The bound for \(\log d_E\) is \(|C_A|z\), and the bound
for \(\log M_0\) is \(18L+2\).

Let \(U\ge2\), count exponents of \(U\) by addition
under multiplication, and charge a sum of two terms by
one plus the larger exponent. Since \(\operatorname{rd}(E)\le d_E\)
and \(\log v\le v\) for \(v>0\), this gives
\(N_0\le U^3\), 
\[
  H_0\le U^{13},
\]
 \(K_0'\le U^{20}\),
\(R_0\le U^{30}\), and \(e_{\rm sp}^{-1}\le U^{50}\).
Also \(\rho_{\rm sp}^{-1}\le U^8\).
Substitution gives \(n',B_{\rm rad}\le U^{100}\), including
the rounding in \(n'\). Thus \(\log n'\le100K_*z\),
while
\[
 \log\Lambda\le(2q_j-1)\max(0,\log B_{\rm rad})
                         \le 200rK_*z .
\]
Because \(K_*z\ge r\) and \(K_*z\ge1\), both are at most
\(200(K_*z)^2\).

Combining them with the coefficient and degree bounds
shows that \(C_{\rm arch}\), \(B_{\rm deg}\), and the
normalized mean of the finite baseline each cost at most
\((17+2K_*+200K_*^4)z^4 S_\kappa\).
It is the upper bound on \(\log\max(1,c_1)\), together
with the fixed frequency restriction, that is used here.
\label{cor:slice-integral-norm-4}

\end{proof}

\paragraph{The lower norm of an integral form of arbitrary weight.}
\label{subsec:deep-cm-interior}
In this paragraph the ambient abelian dimension is \(fr\).
Keep the coefficient-field and level data of the run fixed,
and write \(Z=fr+L+\log d_E+|\Omega|\) in the normalization
of the preceding CM estimates.

\begin{proposition}[Lower norm at the actual weight]\label{prop:slice-lower-integral-norm}
Use the coefficient field and the two-prime level \(g_0g_1\)
of an orientation in Proposition~\ref{prop:linked-comparison-rows}.
Set \(Z=fr+\log b+\log|\operatorname{disc} E|+|\Omega|\).
For a nonzero integral cusp form \(F\) on this run of dominant
weight \(v_0\) with \(|v_0|\ge1\), the normalized supremum norms satisfy
\[
 \operatorname{mean}_\sigma\log\|F\|_{\sup,\sigma}
 \ge -3000Z^2(1+\log|v_0|)|v_0| .
\]
The coefficient \(3000\) is independent of the run, the
weight and the form.
\end{proposition}

\begin{proof}
For the lower norm one uses an interior CM frame of the full
abelian dimension \(g=fr\). Let \(F\) be nonzero in the integral
cusp lattice at a dominant weight \(v_0\), and put \(S=|v_0|\ge1\).
The normalized supremum at the distinguished good embedding
is finite and positive. Take a frame where the reading exceeds
half the supremum. The quantitative CM frame selection starts
with this frame. Write \(p_0\le r\) for the positive rank of
its Hermitian space in the signed frame convention and put
\[
 d_E=|\operatorname{disc}E|,\qquad
 M=g_0g_1\,\mathrm N\mathfrak c_K\,d_{\rm lev},\qquad
 e_0=\min\left\{1,\frac{\pi}{M(fS+2)k_0\log 2}\right\}.
\]
Here \(d_{\rm lev}\) is the positive level discriminant in the
transvection parameter and \(k_0>0\) is the polarization scale of
the component lattice. The principal level gives the integral
real transvections with period \(M\).

At a half-maximal frame,
the nonnegative Fourier modes along these transvections
force at least one majorant \(e_0\) at every nonzero lattice
vector in the reference lattice. Indeed, if all the majorants
\(H_\tau\) are less than \(e_0\le1\), the integral Hermitian
norm and the bound of each of its conjugates by \(H_\tau\)
force that norm to be zero. Along the isotropic flow the
domain condition is \(1+\operatorname{Im}\zeta\,k_0 H_\tau>0\).
Set 
\[
  t_0=1/(2k_0e_0).
\]
 The flow section is then holomorphic
on a strip containing \(\operatorname{Im}\zeta=-t_0\),
is \(M\)-periodic, and tends to zero at the cusp.

On
\(\operatorname{Im}\zeta=-t_0\) each of the factors in the
majorant change lies between \(1/2\) and \(1\). The
weight-spread sum is at most \(S\), so the raw section
has norm at most \(2\cdot 2^{S/2}\) times its modulus at
the starting frame. This uses the half-supremum bound
for the normalized section. Under \(u=e^{2\pi i\zeta/M}\)
the function extends holomorphically across \(u=0\)
with value zero. Apply the maximum principle to its
quotient by \(u\) in the disc of radius \(e^{2\pi t_0/M}\).
At \(\zeta=0\) it yields

\[
  1\le 2^{1+S/2}\exp(-2\pi t_0/M).
\]
The definition of \(e_0\) gives
\[
 2\pi t_0/M\ge(fS+2)\log 2>(1+S/2)\log 2,
\]
a contradiction.

\proofstep{Moving to an interior frame.}
Orthonormalize the starting flag with its order and signs
unchanged. The normalized reading and the positive majorants
are preserved. Put
\[
 \rho_0=(10\sqrt{rS})^{-1},\qquad \rho=3\rho_0/5,\qquad
 \mu=4(fr)^3d_E\sqrt{f/e_0},\qquad A=4\sqrt r\,\mu/\rho.
\]
The majorant lower bound implies the hypothesis of the previous
lattice estimate with \(\lambda=\sqrt{e_0/f}\). Thus \(\mu\)
bounds both the covering radius and the lengths of an independent
spanning family. Approximate the first \(p_0\) axes, scaled by \(A\),
simultaneously at all embeddings by lattice vectors. A matrix with
these vectors as columns has top block \(A I+E_+\) and bottom
block \(E_-\), with \(\|E_\pm\|\le\sqrt r\,\mu\le A\rho/4\).

The columns of this matrix span the positive graph of 
\[
  Z=E_-(A I+E_+)^{-1},
\]
 and
\(\|Z\|\le\rho/3=\rho_0/5\) at every embedding. The absence
of a positive block gives the zero graph. Taking the graph as
the positive part and its Hermitian orthogonal as the negative
part defines the common CM structure on these lattice subspaces.

For the complementary letters let \(G\) be the Gram matrix
over \(E\) of the positive columns. It is positive definite
at every inducing embedding. The orthogonal projection onto
the positive space is given by \(G^{-1}\) in the coordinates
of those columns. Multiplying the complementary projection
by \(\det G\) therefore carries the short lattice vectors into
the original lattice by the adjugate formula and integral
Hermitian pairings.
The matrix identity clearing the projection is
\[
  (\det G)G^{-1}=\operatorname{adj}(G).
\]
 The resulting vectors span the orthogonal
complement. Select an independent family there.

At each
embedding the graph has norm less than \(1/3\), and the
complementary projection in the original Euclidean norm
has norm at most \((1+\|Z\|^2)/(1-\|Z\|^2)\le2\).
This follows by writing the negative graph as \((Z^*y,y)\)
and inverting its Gram matrix \(-(I-ZZ^*)\) for the projection.
Hadamard on the positive Gram matrix gives
\(|\tau(\det G)|\le(2A)^{2p_0}\). Hence the positive
and negative letters all have absolute Hermitian lengths
at most \((2A)^{2p_0+1}\). In the absence of one block
the bound follows directly for the other block by the
same approximation or short basis. Substitution of the approximation radius gives
\[
 \Lambda=X^{2p_0+1},\qquad
 X=\frac{1600}{3}\,r\sqrt f\sqrt S\,(fr)^3 d_E\,e_0^{-1/2}.
\]
The upper estimate on \(\Lambda\) uses any larger bound
for the chosen letters as well.

\proofstep{Retaining a nonzero value through the move.}
To see that the new frame still detects \(F\), use the
holomorphic chart with positive graph \(tZ\) and unchanged
negative coordinate axes. For each signed leading Gram determinant \(\Delta_i\) in this frame, the chart gives
\[
  (1-\|tZ\|^2)^r\le|\Delta_i|\le1.
\]
For leading blocks confined to the positive axes this is the
Gram matrix \(I-(tZ)^*(tZ)\) restricted to those axes.

When negative axes are included, the Schur complement removes
their projection from \(tZ\), and gives the same bound.
The full determinant has absolute value one. Thus for
\(|t|<5\) the change in the log norm factor, with weight
spread at most \(S\), has absolute value at most
\(4Sr\|tZ\|^2\le1/25\), by summing the log minor bounds.
The half-maximal hypothesis bounds the raw reading on
that disc by \(2e^{1/25}<3\) times its modulus at zero.
If the raw reading \(F_0(t)\) vanished at \(t=1\), apply
Schwarz's lemma to

\[
  H(u)=F_0(1+4u)/(3|F_0(0)|)
\]
 on the unit disc.
We have \(H(0)=0\), so
\[
      \tfrac13=|H(-1/4)|\le\tfrac14 .
\]
This contradiction proves that the moved frame still
detects the section.
The graph's majorant is within a factor
\((1+\|Z\|)/(1-\|Z\|)\le2\) of the original. Indeed the
matrix with the positive and negative graph columns is
\(I+J_Z\), where \(J_Z\) is self-adjoint with off-diagonal
blocks \(Z,Z^*\). Consequently the new majorant, represented by \(H_{\rm new}\) in the old frame, satisfies
\[
  H_{\rm new}=(I-J_Z)(I+J_Z)^{-1}.
\]
 This proves the
majorant comparison without changing the polarization.

\proofstep{Polynomial evaluation on the algebraic CM model.}
The arithmetic CM realization uses the lattice of this
basis, its two CM types and the original marking. The
tangent coordinates and the normalized theta generator
are algebraic over a finite extension \(K'\); all finite
comparisons use the corresponding integral models.
\label{cor:slice-integral-norm-5}

The section at this CM point gives a homogeneous polynomial
\(G(X')\) of frame degree \(\mathcal A\), in the entries of the
block matrices \(X'\).
With \(X'\) the frame variables, homogeneity means
\[
  G(cX')=c^{\mathcal A}G(X') .
\]
 Since the section is nonzero in this frame, \(G\ne0\). Choose a nonzero coefficient \(a_\mu\).
At every finite place normalize the theta line and use an
integral eigenbasis of the Lie lattice. The defining
integrality test on \(F\) bounds the normalized polynomial
by one on integral changes of that basis with unit block
determinants. Reduction over an extension with infinite
residue field, or over arbitrarily large finite residue
extensions, gives integral coefficients by the same
determinant test as for the arc. The \(I(\ell_k)\)-scaled
CM letters are integral in this basis. Substituting their
integral coordinate matrices preserves this integrality.
Undoing the scaling gives
\(a_\mu\in I(\mathcal L_\mu)\) at every place, with

\[
  \mathcal L_\mu=(\mathcal T^\vee)^n
\otimes\bigotimes_{\tau,i,k}(\ell_{\tau k}^\vee)^{\mu_{\tau i k}}.
\]
The finite order defect of the toroidal arc is absent in
this interior calculation.
\label{cor:slice-integral-norm-6}

\proofstep{Comparing the archimedean and integral coefficient bounds.}
At \(\sigma\) take the polydisc of radii
\(q_\sigma(\ell_{\tau k})^{-1}\) in the CM matrices.
Each resulting row has length at most \(r\Lambda\) in the
Riemann metric, by the triangle inequality. For invertible
blocks, orthogonalize the rows in the ordering of the flag.
The triangular factor acts by the weight character.
Each diagonal length of that factor is at most the
length of the original row, and all its flag exponents
are nonnegative.

In the orthonormal frame the supremum
bound applies with the fixed theta metric. Consequently
\[
 |G(X')|q_\sigma(\mathcal T)^{-n}
    \le \|F\|_{\sup,\sigma|K}(r\Lambda)^{\mathcal A}.
\]
The bound extends to singular blocks by continuity.
Cauchy's formula in the entries of \(X'\) then bounds
\(|\sigma a_\mu|q_\sigma(\mathcal L_\mu)\) by the same
right side. The coefficient is integral in its line, so
\[
 -\operatorname{ndeg}\mathcal L_\mu
       \le \operatorname{mean}_\sigma
                      \log(|\sigma a_\mu|q_\sigma(\mathcal L_\mu)).
\]
This is the product formula for \(a_\mu\), using
\(\operatorname{ord}_w a_\mu\ge\operatorname{ord}_w I(\mathcal L_\mu)\).
Apply \eqref{eq:slice-cm-line-degree} with \(g=fr\) and
cancel the multiplicity of extensions of each embedding of
\(K\). It follows that
\begin{equation}\label{eq:slice-cm-point-norm}
  \operatorname{mean}_\sigma\log\|F\|_{\sup,\sigma}
       \ge -P_\Lambda(\mathcal A+frn)
                   -\mathcal A\log(r\Lambda).
\end{equation}
This bound is independent of the degree of the CM field over which the chosen data are defined.

\proofstep{Summing the dependence on the weight.}
For the norm budget set \(Z=fr+L+\log d_E+|\Omega|\).
The chosen principal level gives
\(\log(Mk_0)\le 19L+3+\log d_E\), using the two level
prime bounds and the upper bound for \(k_0\).
Since \((fS+2)\log2\le 3fS\) and \(\pi>3\),
\[
 -\log e_0\le 19L+6+\log d_E+f+\log S .
\]

Taking the logarithm of \(X\) gives
\(0\le\log X\le600 Z+\log S\), using
\(\log(1600/3)<7\), \(Z\ge fr\ge1\) and
\(\log fr\le fr\). Thus
\(\log\Lambda\le1800 Z^2(1+\log S)\).
The weight convention gives
\(\mathcal A+frn\le S\) and \(\mathcal A\le S\).
Inserting these bounds into \eqref{eq:slice-cm-point-norm}
yields
\begin{equation}\label{eq:slice-integral-norm}
        \operatorname{mean}_\sigma\log\|F\|_{\sup,\sigma}
                    \ge -3000 Z^2(1+\log |v_0|)|v_0|.
\end{equation}
The absolute coefficient \(3000\) works on both runs,
also when \(v_0\) is the large multiple of the light weight
used in the Chow estimate. The estimates for \(X\) and the
letter metric are the necessary dependence on the weight
of that multiple. They permit the choice of the nonzero
coefficient to vary with \(F\).
\label{cor:slice-integral-norm-7}
\end{proof}
